\begin{tabularx}{\textwidth}{L{37mm}X X}
\toprule
\textbf{Lépés} & \textbf{Feladat} & \textbf{Tipikus eredmény} \\
\midrule
Szakterület megismerése & Az elemzők megértik az alkalmazási terület fogalmait, szabályait, szereplőit és jelenlegi folyamatait. & Szakterületi szójegyzék, háttérismeret, kezdeti modellek. \\
Követelmények összegyűjtése & Együttműködés a kulcsfigurákkal, interjúk, megfigyelés, dokumentumelemzés, forgatókönyvek és prototípusok használata. & Nyers, részben strukturálatlan igénylista és problémafelvetések. \\
Osztályozás és rendezés & A követelmények csoportosítása funkció, nézőpont, prioritás, alrendszer vagy követelménytípus szerint. & Áttekinthető követelménycsoportok, kapcsolatok, duplikációk felismerése. \\
Ellentmondások feloldása & Az eltérő érdekű kulcsfigurák igényeinek egyeztetése, konfliktusok feltárása és kezelése. & Egyeztetett döntések, módosított vagy elvetett követelmények. \\
Fontossági sorrend felállítása & A követelmények üzleti érték, kockázat, költség és sürgősség szerinti rangsorolása. & Prioritási lista, inkremensek vagy kiadási terv alapja. \\
Követelményellenőrzés & Annak vizsgálata, hogy a követelmények teljesek, konzisztensek, megvalósíthatók, érthetők és a valódi igényeket tükrözik-e. & Javított követelményhalmaz, validációs megjegyzések, nyitott kérdések. \\
\bottomrule
\end{tabularx}
